本章对 $J$-全纯曲线作整体研究。首先给出 $J$-全纯曲线基本模空间的定义，然后在 $C^\infty$ 情形下给出孔采维奇 (Kon95) 引入的稳定映射模空间之定义。特别地，我们紧随深谷贤治与小野薫 (FOn99) 给出的证明，说明\emph{稳定映射拓扑}的精确定义，并证明其紧性与豪斯多夫性。

我们将大体局限于闭情形；关于带边黎曼曲面的完整讨论，请读者参阅 (FOOO09)、(L02)。这里只引用其中研究本书第二卷拉格朗日弗洛尔同调所需的基本定理。

\section{伪全纯曲线的模空间}
\stfhIndex{I205}{mo^kong^jian}{模空间（moduli space）}
\label{sec:moduli-space-pseudoholomorphic-curves}

格罗莫夫的非线性柯西--黎曼方程，是关于“紧定向曲面 $\Sigma$ 的共形结构以及从 $\Sigma$ 到殆凯勒流形 $(M,\omega,J)$ 的映射”的方程。本节给出这种偶对 $(j,u)$ 之模空间的定义。

先描述基本研究对象。固定一个欧拉示性数为
\[
 \chi(\Sigma)=2-2g
\]
的闭定向曲面 $\Sigma$，其中 $g$ 是其亏格，$2g=\operatorname{rank}H_1(\Sigma,\mathbb Z)$。以 $[\Sigma]$ 表示 $\Sigma$ 的基本类。给定连续映射 $u:\Sigma\to M$ 时，它定义自然元素
\[
 u_*[\Sigma]\in H_2(M,\mathbb Z),
\]
记为 $[u]$。以 $A$ 表示 $H_2(M,\mathbb Z)$ 的一般元素。当 $u$ 可微时，积分
\[
 \int u^*\omega
\]
有定义，并且只依赖于同调类 $[u]=u_*[\Sigma]$。它称为\emph{可微}映射 $u$ 的\emph{辛面积}。
\stfhIndex{I270}{xin^mian^ji}{辛面积（symplectic area）}

现在固定曲面的一种微分同胚型，并以 $\Sigma$ 表示任意这种曲面。考虑由“$\Sigma$ 的共形结构与一个映射”构成的偶对 $(j,u)$。注意，群 $\Diff(\Sigma)$ 通过推前作用于偶对集合：
\[
 (\phi,(j,u))\longmapsto(\phi_*j,u\circ\phi^{-1}).
\]
记
\[
 \overline\partial_J(j,u)=\frac{du+J\circ du\circ j}{2},
\]
以强调它对定义域复结构 $j$ 的依赖。在 $\Sigma$ 上固定 $j$ 的一个埃尔米特度量 $h$，并在 $M$ 上固定任意参考相容度量 $g$。

我们希望研究以 $(j,u)$ 为变量的柯西--黎曼方程之解集：
\[
 \overline\partial_J(j,u)=0,
 \qquad \int|du|^2<\infty.
\]
这里，能量密度 $|du|^2$ 用度量 $h$ 和 $g$ 计算。
\stfhIndex{I115}{neng^liang^mi^du}{能量密度（energy density）}

对具有非空边界 $\partial\Sigma$ 的带边黎曼曲面 $\Sigma$，考虑满足
\[
 \overline\partial_J(j,w)=0,
 \qquad w(\partial\Sigma)\subset L,
 \qquad \int|dw|^2<\infty
\]
的映射 $w:\Sigma\to M$，其中 $L$ 是给定的拉格朗日子流形。

就边值问题的椭圆性而言，可以用全实子流形替换拉格朗日子流形 $L$，正如上一章所说明的。下面的引理是拉格朗日子流形区别于全实子流形的关键几何性质。

\begin{lemma}\label{lem:lagrangian-relative-homology-symplectic-area}
设 $L$ 是 $(M,\omega)$ 的拉格朗日子流形。假设两个映射
$w_1,w_2:(\Sigma,\partial\Sigma)\to(M,L)$ 具有相同的相对同调类
$[w_1]=[w_2]\in H_2(M,L)$。则
\[
 \int w_1^*\omega=\int w_2^*\omega.
\]
\end{lemma}

\begin{proof}
由假设，链 $w_1-w_2$ 同调于支集包含在 $L$ 中的链 $C$；也就是说，$M$ 中存在一个 $3$-链 $D$，使得
\[
 \partial D=w_1-w_2-C.
\]
由于 $d\omega=0$ 且 $\int_C\omega=0$，斯托克斯公式蕴含
\[
 0=\int_{\partial D}\omega
 =\int w_1^*\omega-\int w_2^*\omega-\int_C\omega
 =\int w_1^*\omega-\int w_2^*\omega,
\]
这就完成了证明。
\end{proof}

所以，如果把映射 $u$ 的同调类固定为
\[
 u_*[\Sigma]=A\in H_2(M),
\]
或者在拉格朗日边界 $L$ 的情形固定为
\[
 w_*[\Sigma,\partial\Sigma]=B\in H_2(M,L),
\]
则映射 $u$（相应地，$w$）的辛面积保持不变。于是，在这两种情形的任一种中，对任意 $J$-全纯映射 $u$，都有恒等式
\begin{equation}
 \int u^*\omega=E_J(u)=\frac12\int|du|^2,
 \label{eq:j-holomorphic-symplectic-area-energy}
\end{equation}
其中以 $E_J(u)$ 表示 $u$ 关于度量 $g_J=\omega(\cdot,J\cdot)$ 的调和能量。

不过在实践中，需要让殆复结构在 $\mathcal J_\omega$ 中变化，或者考虑依赖于 $u$ 的定义域变量的一族这种殆复结构。因此在实际几何估计中，需要固定参考度量，并把 $J$ 本身视作辅助参数。为此，下面这个先验能量估计的参数化版本非常重要。

\begin{lemma}\label{lem:uniform-energy-varying-j}
设 $K\subset\mathcal J_\omega$ 是紧子集，并假设 $M$ 紧。在 $M$ 上固定任意度量 $g$。设 $(j,u)$ 对某个 $J\in K$ 是 $J$-全纯映射，并以 $E(u)$ 表示 $u$ 关于度量 $g$ 的调和能量。假设 $[u]=A\in H_2(M)$（相应地，$[u,\partial u]=B\in H_2(M,L)$）。则存在常数 $C=C(K)$，使得对所有 $J\in K$ 的 $J$-全纯映射 $(j,u)$ 都有
\[
 E(u)\leq C\omega(A)
\]
（相应地，$E(u)\leq C\omega(B)$）。此外，关于 $\mathcal J_\omega$ 的 $C^0$ 拓扑，常数 $C(K)$ 在 $K$ 上连续变化。
\end{lemma}

\begin{proof}
对 $J\in K$，比较 $E_J(u)$ 与 $E(u)$。回忆度量在逐点意义下是正定二次型，并且存在连续的常数 $c(x),C(x)>0$，使每个 $x\in M$ 上都有
\[
 c(x)g(x)<g_J(x)<C(x)g(x).
\]
记
\[
 c_J=\min_{x\in M}c(x),
 \qquad C_J=\max_{x\in M}C(x).
\]
$c_J,C_J$ 关于 $J$（在 $C^\infty$ 拓扑中）的连续性是立即的。于是
\[
 c_J|du|^2\leq|du|_J^2\leq C_J|du|^2,
\]
从而 $c_JE(u)\leq E_J(u)\leq C_JE(u)$。由于 $E_J(u)=\omega(A)$，得到 $E(u)\leq(1/c_J)\omega(A)$。令
$C=C(K)=\min_{J\in K}1/c_J$，便得到该不等式。常数 $C(K)$ 在 $K$ 上的连续性立即来自表达式
$C(K)=\min_{J\in K}1/c_J$。
\end{proof}

现在研究解的模空间
\begin{equation}
 \widetilde{\mathcal M}^{E_0}(\Sigma,J)
 =\left\{(j,u)\mathrel{\Big|}
 \begin{array}{l}
 u:(\Sigma,j)\to(M,J),\quad \overline\partial_{(J,j)}u=0,\\
 \displaystyle\int|du|_J^2\leq E_0,\quad u(\partial\Sigma)\subset R
 \end{array}\right\}
 \label{eq:bounded-energy-moduli-space}
\end{equation}
及其（在 $C^\infty$ 拓扑中的）收敛性质，其中\emph{给定常数 $E_0>0$}。

\begin{remark}\label{rem:finite-energy-automatic-fixed-class}
由上述引理，如果把同调类 $A$ 固定为
\[
 u_*[\Sigma]=A\in
 \begin{cases}
 H_2(M),&\partial\Sigma=\varnothing,\\
 H_2(M,L),&\partial\Sigma\ne\varnothing,
 \end{cases}
\]
那么在 $\partial\Sigma=\varnothing$ 或 $R$ 是拉格朗日子流形 $L$ 的情形，有限能量条件将自动成立。
\end{remark}

设 $\{(j_\alpha,u_\alpha)\}$ 是 $\widetilde{\mathcal M}^{E_0}(\Sigma,J)$ 中的序列。收敛失败有如下几个显然的来源：
\begin{enumerate}[label=(\arabic*)]
 \item $\max_{z\in\Sigma}|du_\alpha(z)|\to\infty$；
 \item 当 $g(\Sigma)\ne0$ 时，$j_\alpha$ 可能趋向 $\Sigma$ 的共形结构空间之无穷远处，即可能退化；
 \item 当 $g(\Sigma)=0$ 时，重参数化群 $PSL(2,\mathbb C)$（对 $S^2$）或 $PSL(2,\mathbb R)$（对 $D^2$）非紧；也就是说，如果
 $(j,u)\in\widetilde{\mathcal M}^{E_0}(\Sigma,J)$，则对任意
 $\phi\in PSL(2,\mathbb C)$ 或 $PSL(2,\mathbb R)$，都有
 $(\phi^*j,u\circ\phi)\in\widetilde{\mathcal M}^{E_0}(\Sigma,J)$。
\end{enumerate}

在 $g(\Sigma)=0$ 的情形，通过研究商空间
\[
 \mathcal M^{E_0}(\Sigma,J)
 =\widetilde{\mathcal M}^{E_0}(\Sigma,J)/\sim
\]
上关于商拓扑的收敛，消除 $\widetilde{\mathcal M}^{E_0}(\Sigma,J)$ 因重参数化而产生的非紧性；该商空间是等价类 $[u_\alpha]$ 的集合。

换言之，$[u_\alpha]$ 收敛是指存在 $\phi_\alpha\in\operatorname{Aut}(\Sigma)$，使 $u_\alpha\circ\phi_\alpha$ 收敛。研究模空间紧性时，一个困难来源是没有选择重参数化 $\phi_\alpha$ 的典范方式；这要求在研究收敛时根据具体情形作一定程度的巧妙处理。
